\begin{table*}[t]
\centering
\caption{
\textbf{Order-Sensitivity Diagnostics and Dependency-Violation Sensitivity, CaptainCookRecipes (GPT-4o-mini).}
TopoVar is the variance-based companion to TopoDev (Table~\ref{tab:valid_reordering}):
mean per-step sample variance across sampled orderings, more robust to a single outlier
ordering than the max-min range. TopoVar is on a \emph{squared}-score scale, so it is not
directly comparable to a raw score or threshold; $\sqrt{\text{TopoVar}}$ (a standard
deviation) brings it back to the original $[0,1]$ score scale and is the quantity to
compare against the method's own error-detection threshold (also on that scale).
Notably, ARES's typical reordering-induced swing ($\sqrt{\text{TopoVar}}=0.156$) is
nearly equal to its own decision threshold (0.175)---ordinary, harmless reordering
noise alone is large enough to plausibly flip a claim across the sound/unsound boundary.
Error-Recovery Jaccard is the mean pairwise Jaccard
similarity between the sets of nodes flagged as errors under two different valid
orderings of the same recipe (1.0 = predicted error set never changes with ordering).
Exact-Match Rate is the fraction of orderings whose predicted error set exactly equals
the ground-truth error set. Dependency-Violation Correct-Direction Rate (Experiment~2) is
the fraction of minimal single-edge dependency-violating swaps for which the score
actually \emph{decreased}, with a two-sided binomial test against a 50\% chance baseline;
$N=245$ violation cases across 24 recipes. Notably, this rate is \emph{hyperparameter-
sensitive for ARES}: at the faster $\epsilon=0.3$ setting used during development, ARES's
correct-direction rate was 34.7\% ($p=2\times10^{-6}$, significantly \emph{worse} than
chance); at the paper-faithful $\epsilon=0.1$ shown here it is 58.8\%
($p=0.007$, significantly \emph{better} than chance) --- the two epsilon settings give
opposite-sign conclusions. Entail-Prev and Entail-Base do not use $\epsilon$ and are
stable across both settings (53.1\%/52.2\% and 6.1\%/8.2\% respectively), consistent with
the ARES-specific shift being a real epsilon effect rather than run-to-run noise.
\emph{Mechanism:} ARES's \texttt{cert\_nonexact} scorer estimates each step's soundness
by averaging over $N\!\approx\!\log(2m/\delta)/(2\epsilon^2)$ randomly sampled premise
subsets ($\sim$265 samples/step at $\epsilon=0.1$ vs.\ $\sim$30 at $\epsilon=0.3$); the
coarser $\epsilon=0.3$ estimate is noisy enough to invert the \emph{aggregate direction}
of the measured effect, not merely blur its magnitude. Even at $\epsilon=0.1$ the
correct-direction rate is only $\sim$9 points above chance, and the effect exceeds
ordinary valid-reordering noise (Exceeds-Noise rate, not shown) in only 15.1\% of cases
--- a real but weak signal, not a reliable detector. This hyperparameter sensitivity is
specific to the dependency-violation test; it does not affect the valid-reordering
findings above (TopoDev/TopoVar/Error-Recovery replicate at both $\epsilon$ settings),
which remain this paper's central claim.
All columns use $\epsilon=0.1$ (paper-faithful, $N=24$ recipes).
}%
\label{tab:diagnostics}
\resizebox{\textwidth}{!}{
\begin{tabular}{lccccc}
\toprule
\textbf{Method}
& \textbf{TopoVar $\downarrow^{\ddagger}$ ($\sqrt{\text{TopoVar}}$)}
& \textbf{Threshold$^{\ddagger}$}
& \textbf{Error-Recovery Jaccard $\uparrow^{\ddagger}$}
& \textbf{Exact-Match Rate $\uparrow^{\ddagger}$}
& \textbf{Dep.-Violation Correct-Direction Rate} \\
\midrule

ARES
    & 0.024 \small{[0.019, 0.030]} (0.156 \small{[0.137, 0.174]})
    & 0.175
    & 0.734
    & 0.4\%
    & 58.8\% ($p = 0.007$) \\
Entail-Prev
    & 0.081 \small{[0.068, 0.093]} (0.284 \small{[0.260, 0.305]})
    & 0.400
    & 0.604
    & 0.0\%
    & 53.1\% ($p = 0.37$, n.s.) \\
Entail-Base
    & 0.032 \small{[0.023, 0.041]} (0.179 \small{[0.153, 0.202]})
    & 1.000
    & 0.903
    & 0.0\%
    & 6.1\% ($p < 10^{-4}$) \\

\bottomrule
\end{tabular}
}
\end{table*}
